% !TEX program = xelatex
\documentclass{英语试卷}
\begin{document}
\examtitle{重庆市高2025届高三第三次质量检测}{英语}
\par 重庆市高2025届高三第三次质量检测\par
\par 高三英语 试题卷\par
\par 命审单位：重庆南开中学\par
\begin{notice}
\par 1. 答卷前，考生务必将自己的姓名、准考证号填写在答题卡上。\par
\par 2. 回答选择题时，选出每小题答案后，用铅笔把答题卡上对应题目的答案标号涂黑。如需改动，用橡皮擦干净后，再选涂其他答案标号。回答非选择题时，将答案写在答题卡上，写在本试卷上无效。\par
\par 3. 考试结束后，将本试卷和答题卡一并交回。微信公众号 IAI English\par
\end{notice}
\parttitle{第一部分 听力（共两节，满分30分）}
\par 做题时，先将答案标在试卷上，录音内容结束后，你将有两分钟的时间将试卷上的答案转涂到答题卡上。微信公众号IAI English\par
\sectitle{第一节（共5小题：每小题1.5分，满分7.5分）微信公众号IAI English}
\par 听下面5段对话。每段对话后有一个小题，从题中所给的A、B、C三个选项中选出最佳选项。听完每段对话后，你都有10秒钟的时间来回答有关小题和阅读下一小题。每段对话仅读一遍。\par
\question{1. What will the man do first?}
\begin{choices}(3)
  \item Have a meal.
  \item Write a report.
  \item Meet a client.
\end{choices}
\question{2. When does the woman say she'll arrive?}
\begin{choices}(3)
  \item At 2:10.
  \item At 2:30.
  \item At 2:50.
\end{choices}
\question{3. What is the probable relationship between the speakers?}
\begin{choices}(3)
  \item Husband and wife.
  \item Salesman and customer.
  \item Employer and employee.
\end{choices}
\question{4. Why does the woman speak to the man}
\begin{choices}(3)
  \item To ask for information.
  \item To offer an invitation.
  \item To make an apology.
\end{choices}
\question{5. What are the speakers probably talking about?}
\begin{choices}(3)
  \item A train schedule.
  \item A natural disaster.
  \item A travel experience.
\end{choices}
\sectitle{第二节（共15小题：每小题1.5分，满分22.5分）}
\par 听下面5段对话或独白。每段对话或独白后有几个小题，从题中所给的A、B、C三个选项中选出最佳选项。听每段对话或独白前，你将有时间阅读各个小题，每小题5秒钟；听完后，各小题将给出5秒钟的作答时间。每段对话或独白读两遍。微信公众号IAI English\par
\par 听第6段材料，回答第6至7题。\par
\question{6. What do we know about the woman's ring?}
\begin{choices}(3)
  \item She made it herself.
  \item She bought it at a market.
  \item She received it as a present.
\end{choices}
\question{7. What did the woman enjoy about the experience?}
\begin{choices}(3)
  \item She made some friends.
  \item She had some quiet time.
  \item She saw some great artwork.
\end{choices}
\par 听第7段材料，回答第8至9题。\par
\question{8. How does the man feel about riding in a driverless car?}
\begin{choices}(3)
  \item Bored.
  \item Excited.
  \item Stressed.
\end{choices}
\question{9. How many concerns does the woman mention about driverless cars?}
\begin{choices}(3)
  \item Two.
  \item Three.
  \item Four.
\end{choices}
\par 听第8段材料，回答第10至12题。\par
\question{10. Where are the speakers most probably?}
\begin{choices}(3)
  \item In a restaurant.
  \item In a stadium.
  \item In a library.
\end{choices}
\question{11. How does the woman sound hearing the man's first suggestion?}
\begin{choices}(3)
  \item Pleased.
  \item Annoyed.
  \item Surprised.
\end{choices}
\question{12. What will the woman help the man do?}
\begin{choices}(3)
  \item Organize the competition.
  \item Give the swimming class.
  \item Invite the vice-president:
\end{choices}
\par 听第9段材料，回答第13至16题。\par
\question{13. What was the woman's first job?}
\begin{choices}(3)
  \item A club DJ.
  \item A band manager.
  \item A marketing assistant.
\end{choices}
\par 14 Why did the woman leave her previous job? A. To further her education. B. To take care of her brother. C. To look for a more exciting job.\par
\question{15. What did Eric Johnson agree to let the woman do?}
\begin{choices}(3)
  \item Sit in on his show.
  \item Take the trainee course.
  \item Do weather reports on air.
\end{choices}
\question{16. How many days a week did the woman work as a waitress?}
\begin{choices}(3)
  \item Three days.
  \item Four days.
  \item Five days.
\end{choices}
\par 听第 10 段材料，回答第 17 至 20 题。\par
\question{17. Why did the speaker start to take an interest in gorillas?}
\begin{choices}(3)
  \item Because of a film.
  \item Because of a book.
  \item Because of a poster.
\end{choices}
\question{18. What did the speakers uncle do for him?}
\begin{choices}(3)
  \item He took him on a trip.
  \item He gave him many photos.
  \item He taught him how to read.
\end{choices}
\question{19. What surprised the speaker about gorillas?}
\begin{choices}(3)
  \item They have no tails.
  \item They barely drink water.
  \item They eat insects sometimes.
\end{choices}
\question{20. What will the speaker introduce next about gorillas?}
\begin{choices}(3)
  \item Their main habitats.
  \item Their social behavior.
  \item Their protection issues.
\end{choices}
\parttitle{第二部分 阅读理解（共两节，满分50分）}
\sectitle{第一节（共15小题：每小题2.5分，满分37.5分）}
\par 阅读下列短文，从每题所给的A、B、C和D四个选项中，选出最佳答案。\par
\passagelabel{A}
\par The American's Favorite Teacher competition saw many hopeful educators competing for\$25,000, a trip to Hawaii, and a professional photo shoot, all to raise awareness for Teach For America. At the final bell, the finalists emerge victorious.\par
\par Kelby Milgrim\par
\par Mr. Milgrim, or Sir, to his students, is the STEM director at Phoenix Christian Preparatory School and helped establish the Phoenix chapter of the Southwest National Robotic Challenge. By exposing students to diverse experiences in science, technology, engineering and mathematics, he prepares them for future challenges and opportunities. Kelby also teaches robotics, honors physics and forensics.\par
\par Shelly Starline\par
\par Mrs. Starline teaches middle school language arts in Cincinnati. She is an empathetic educator passionate about the transformative power of literacy. Shelly brings the fun into the classroom by incorporating new activities, unusual facts, and interesting current events to foster creativity and critical thinking. She is grateful for the support of her family, especially her children and husband, Jeremy.\par
\par Erica Boyce\par
\par Coach, Beyonce, Ms. B or Mom, as her students call her at Bronx Theater High School and YABC, is a special education generalist from Brooklyn. "Teaching is less of a passion and more of my purpose," she said. Erica is also the proud coach of the school's tennis team. Realizing sports' impact on teens, she offers free tennis lessons to kids and welcomes support.\par
\par Jasmine Williams\par
\par Ms. Jasmine is a lively teacher in Chicago's West Loop area. Her passion for early childhood development is evident as she brings excitement into the classroom through songs, dance, and art. Jasmine loves her family and would like to give a shoutout to her mother Angela Williams and sister Charlene Rhinehart for always encouraging her to pursue her dreams and for speaking positive words over her.\par
\question{21. How many rewards will America's favorite teacher get?}
\begin{choices}(2)
  \item Two.
  \item Three.
  \item Four.
  \item Five.
\end{choices}
\question{22. Which teacher is most suitable for you if you have a 3-year-old daughter?}
\begin{choices}(2)
  \item Kelby Milgrim.
  \item Shelly Starline.
  \item Erica Boyce.
  \item Jasmine Williams.
\end{choices}
\question{23. In which section of a magazine will this passage appear?}
\begin{choices}(2)
  \item Science.
  \item Fiction.
  \item Education.
  \item Economy.
\end{choices}
\passagelabel{B}
\par A graduate from the renowned Colorado State University Veterinary Program, Dr Kwane Stewart has been a California vet (兽医) for more than two decades.\par
\par Many years ago during a five-year service as a shelter vet, Stewart thought about quitting. He’d long dreamed of saving animals but was instead forced to euthanize (使……安乐死) an increasing number of animals. "I’m seeing lots of unwanted pets dropped off by people who don’t have the money to feed them or care for them medically, which started to steal a part of my soul. I thought about leaving the veterinary profession altogether." It was also during the service that he began to form a true understanding of those that struggle to afford health care for their pets.\par
\par Project Street Vet came about after his encounter with a homeless man and his dog outside the 7-11. Stewart treated the dog's skin condition and the animal was transformed. But for Stewart, the man's gratitude was a wake-up call. "Thank you for not ignoring me" were the words that inspired his next chapter.\par
\par Soon after, Stewart began to do more outreach, setting up small drop-in clinics to provide medical care to pets whose owners couldn't afford it and walking the streets looking for unhoused individuals whose pets needed help. Right away he was struck by how much the pet owners cared for their beloved animals.\par
\par Stewart found that many pet owners wanted to share their stories about their dog and their history together. With their permission. Stewart's brother eventually began documenting some of the work and personal stories to raise awareness about animal welfare and homelessness.\par
\par His brother encouraged him to join forces with others to provide street care, and soon more like-minded animal welfare professionals jumped on board.\par
\par Stewart's work with Project Street Vet is all volunteer, and the organization has expanded to other cities. "I'm building a network of trusted volunteers, technicians with hospitals and clinics we can call on. These are all people who just reached out and said, 'I'm inspired by what you do. How do I do it?'" Since he started, Stewart and his volunteer teams have treated thousands of animals while giving their human parents hope and dignity, too.\par
\question{24. Why did Kwane Stewart once want to leave his profession?}
\begin{choices}(2)
  \item Because he didn't have a bright career prospect.
  \item Because he saw lots of pets adopted by people.
  \item Because the reality was against his dream.
  \item Because the job was energy-consuming.
\end{choices}
\question{25. What inspired Project Street Vet?}
\begin{choices}(2)
  \item The five-year service as a vet.
  \item The homeless man's gratitude.
  \item The volunteers' involvement.
  \item The dog's recovery.
\end{choices}
\question{26. What do we know about Stewart's brother?}
\begin{choices}(2)
  \item He was a professional vet.
  \item He provided help for Stewart.
  \item He offered to share his pets' stories.
  \item He was encouraged to save animals.
\end{choices}
\question{27. What might be the best title for the passage?}
\begin{choices}(2)
  \item Hold On or Give Up
  \item Save the Deserted Pets
  \item Project Street Vet Needs Help
  \item Kwane Stewart, the Street Vet
\end{choices}
\passagelabel{C}
\par Each Olympic Games features athletes pushing the boundaries of human limits, but the Paris Olympics are particularly notable in this regard. At the Paris Olympics, records were broken through group improvement rather than individual breakthroughs. In the men's 10,000 meters, 13 athletes broke the 2008 Olympic record, with Joshua Cheptegei winning in 26:43.14. The women's 400 meters final saw all nine athletes finish under 50 seconds.\par
\par The amazing purple track may account for that. The purple runway of the Paris Olympics is inspired by lavender fields in Provence, southern France. This color not only symbolizes the romance and elegance of France, but also shows the natural scenery and cultural heritage of France.\par
\par The design concept and purpose of the purple runway is to bring a sense of freshness and a unique viewing experience through non-traditional color selection. Alain Blondel, the head of athletics for the Paris Olympics, commented on the impact of the purple track color. He noted that this design not only creates a strong visual contrast between the track and the spectators during television broadcasts, highlighting the athletes, but also helps athletes perform better. This is due to the significant effect that color can have on an athlete's performance.\par
\par Purple is considered a cool color, which can create a sense of calm and stability. Unlike blue, another cool color, purple also incorporates elements of red, a warm color that adds a sense of energy and expansion. This combination of red's warmth and blue's coolness gives purple a balanced effect— it calms without being too subdued (柔和的), helping to maintain a steady yet motivated emotional state.\par
\par This runway is not only eye-catching for its unique color, but also hides the secrets of high technology. The runway is made of high-tech materials developed by Mondo, which have excellent durability and elasticity, and can provide the best competitive condition. Kipchoge, the marathon world record holder, praised the impact and resilience of the runway. Track and field legend Carl Lewis also conducted a comprehensive test of the runway. He believed that the runway can maintain stable performance in all climate condition and that athletes all over the world would love this track full of creativity and technology.\par
\question{28. How does the author prove his point in the first paragraph?}
\begin{choices}(2)
  \item By making comparisons.
  \item By analyzing causes.
  \item By presenting opinions.
  \item By giving examples.
\end{choices}
\question{29. What does Alain Blondel probably agree with?}
\begin{choices}(2)
  \item Television stations should highlight all the athletes.
  \item Purple has positive effects on an athlete's performance.
  \item The strong visual effect of the track will excite audience.
  \item The purple track strengthens the contact between athletes and the audience.
\end{choices}
\question{30. What advantage does purple have according to the passage?}
\begin{choices}(2)
  \item Creating a sense of calm and stability like blue.
  \item Helping athletes to develop a balanced personality.
  \item Having a stronger sense of energy in comparison with red.
  \item Contributing to keeping a stable yet motivated emotional state.
\end{choices}
\question{31. What is the purpose of mentioning Kipchoge and Lewis in the last paragraph?}
\begin{choices}(2)
  \item To prove the runway can function excellently.
  \item To show they are the legendary figures in the field of sports.
  \item To promote the international influence of the Paris Olympics.
  \item To attract more athletes to test whether the runway is up to standard.
\end{choices}
\passagelabel{D}
\par Orange juice has never been more expensive. Concentrate orange-juice futures in New York have quadrupled (翻两番) since late 2021.\par
\par The problem is that the market for oranges to make concentrate is very concentrated. Growth of the fruit in Florida, once the world's second-biggest producer, has decreased by 92\% in the past 20 years after being hit by pests, hurricanes and rising costs. Other growers, such as Mexico and Spain, largely supply the fresh-fruit market. That leaves Brazil, which accounts for 70\% of the world's production and a similar share of exports. And now it, too, is,in trouble.\par
\par At the start of the year the market already faced bitter prospects. Brazil's drought is getting worse—parts of Sao Paulo, Brazil's main orange-producing state, have seen no rain for four months. In May, growers from Sao Paulo forecast that the year's crop would fall by 24\% compared with the year before, after an ultra-hot flowering season. Farmers also faced an out-of-control resurgence in citrus greening (柑橘黄龙病), an untreatable bacterial disease that turns oranges sour before killing the tree.\par
\par Bigger farmers can afford to protect groves by using nets or insecticides. For smaller ones the only fix is to root out potentially diseased trees. High prices mean many prefer to keep their groves in the ground, so infections persist.\par
\par Planting in other regions may help, but suitable areas lack water access and are too far away from juice plants, which means that progress is very slow. Some farmers may opt to plant less risky crops than oranges, such as sugar cane.\par
\par In many fields damage has already been done, meaning this year's harvest is doomed. The next few could be just as bad. New trees take four years to bear fruit; they too could be exposed to greening. Droughts could sap yields again.\par
\par Despite high prices, global demand for juice remains strong, because consumers are affluent enough to pay more. This may lift producers' profits a little, but not by much, for they also face higher costs.\par
\par Shortages may start to occur soon, probably outside America and Europe, which will be prioritized by retailers as the largest orange-juice markets. Without a miracle, things look rotten for seekers of Vitamin C.\par
\question{32. What contributes to the decline in orange production in Brazil?}
\begin{choices}(2)
  \item Consumer preference towards fresh fruit.
  \item Lack of water and a bacterial disease.
  \item A decrease in the number of farmers.
  \item The introduction of new markets.
\end{choices}
\question{33. What is the impact of citrus greening on orange production?}
\begin{choices}(2)
  \item It increases the sweetness of oranges.
  \item It affects the quality and yield of oranges.
  \item It satisfies the demand for oranges as a health supplement.
  \item It causes a decrease in the number of trees available for planting.
\end{choices}
\question{34. How are same orange farmers likely to address the challenges?}
\begin{choices}(2)
  \item To seek less risky alternatives.
  \item To root out diseased trees by using nets.
  \item To plant oranges far from juice factories.
  \item To turn oranges sour before killing the tree.
\end{choices}
\question{35. What does the underlined word "affluent" in Paragraph 7 probably mean?}
\begin{choices}(2)
  \item Economical.
  \item Influential.
  \item Wealthy.
  \item Considerate.
\end{choices}
\sectitle{第二节（共5小题；每小题2.5分，满分12.5分）}
\par 阅读下面短文，从短文后的选项中选出能填入空白处的最佳选项。选项中有两项为多余选项。\par
\par Wake up early enough and, depending on where you live, you'll likely hear a pleasant chorus of birds chirping (啾啾叫) their tiny hearts out. \blanknum{36}, but why?\par
\par According to the Woodland Trust, this early singing is known as the dawn chorus, and it can start as early as 4 am and last for several hours. \blanknum{37}. One theory as to why they opt for the wee hours to communicate these important messages is because the low visibility makes it hard to do other bird activities, like searching for food. Keeping their activity level low, they opt to sing instead.\par
\par \blanknum{38}. By singing with passion, they're letting other birds know they survived the night—no guarantee in nature—and would make for an excellent mate.\par
\par Like a good recording studio, the early morning hours also allow birds to transmit a clear tune thanks to the cooler, drier air. Because birds have distinctive voices, that clarity helps them be more easily identified by birds within listening distance.\par
\par There used to be a popular theory that birds sang so much in the mornings because the sounds could travel greater distances due to less atmospheric turbulence, but it was disproven by researchers at the University of Western Ontario in 2003. They played recordings of sparrows at dawn and midday. \blanknum{39}, but they were more consistent.\par
\par \blanknum{40}, so singing when they stand the best chance of being heard makes sense. For a bird, an early morning session is like having the perfect acoustic environment(音效环境) for their performance.\par
\begin{choices}(1)
  \item This is natural behavior for birds
  \item The songs didn't travel farther at dawn
  \item Birds only have so much energy to sing one out
  \item Birds do it to attract mates and warn other birds to avoid their turf (地盘)
  \item The early morning singing of the birds is a sign that a new day has begun
  \item Another theory argues that birds use the morning vocals to convey strength
  \item The birds sing at the break of morning in order to create a symphony of nature
\end{choices}
\parttitle{第三部分 语言知识运用（共两节，满分30分）}
\sectitle{第一节（共15小题；每小题1分，满分15分）}
\par 阅读下面短文，从短文后各题所给的A、B、C和D四个选项中，选出可以填入空白处的最佳选项。\par
\par Mo really hated physics. \blanknum{41}, her friend Terry didn't have a problem with physics. Terry had promised to help Mo revise for the physics exam, and then just ignored all her calls and texts. \blanknum{42}, Mo rang her house. but her mum just told Mo she was busy and hung up. "Some good friend!" thought Mo bitterly,\par
\par During the physics exam, she accidentally looked up and saw Terry having her phone on her left knee and reading from it. Is that how Terry always got good grades? It was even harder to \blanknum{43} the exam after that. She was furious with Terry. After a great deal of \blanknum{44}, she decided to tell her teacher Mr. Reed.\par
\par Mo made sure she was sitting behind Terry again in the history exam. They were half way \blanknum{45} when she noticed a movement. She looked up to see Terry with her phone under the table. Mo nodded to Mr. Reed. Terry was so busy with her phone that she didn't even hear him as he \blanknum{46}. He didn't say a word as Terry looked up from her phone with a frightened expression. Outraged, he took her exam paper, \blanknum{47} it up and pointed to the door to tell Terry to leave the room. Everyone stared after her in shock as she stormed out in tears. Mo felt more \blanknum{48} than satisfaction. What had she done?\par
\par Mo was walking towards the school gate when Terry caught up with her. It was \blanknum{49} that Terry had been crying. Her face was pale and her eyes were \blanknum{50}. Terry explained to her that her dad had had a heart attack and had been \blanknum{51} for a week. Terry was so concerned about her dad that she kept \blanknum{52} her mom to see how the operation was going during the exam. Mr. Reed believed her now, but she had to retake the exam. Mo felt so regretful for jumping to conclusions and \blanknum{53} her best friend so easily.\par
\par What we see with eyes is not \blanknum{54} true. When friendship and faith are tested, what matters is to reach out and get the full picture. Real friendship is far too valuable to be thrown away in a moment of emotional \blanknum{55}.\par
\cz{41}{Consequently}{Otherwise}{Contrarily}{Ideally}
\begin{choices}(3)
  \item Frustrated
  \item Motivated C Insulted
  \item Relieved
\end{choices}
\cz{43}{set about}{concentrate on}{prepare for}{participate in}
\cz{44}{anticipation}{consideration}{confrontation}{adaptation}
\cz{45}{in}{out}{through}{up}
\begin{choices}(3)
  \item sneaked up
  \item got in
  \item fell down D cried out
\end{choices}
\cz{47}{tore}{finished}{put}{packed}
\cz{48}{joy}{guilt}{anger}{calmness}
\cz{49}{fortunate}{unlikely}{evident}{funny}
\cz{50}{open}{shut}{sparkling}{swollen}
\cz{51}{cured}{released}{stuck}{hospitalized}
\cz{52}{phoning}{texting}{bothering}{begging}
\cz{53}{trusting}{helping}{judging}{catching}
\cz{54}{occasionally}{generally}{partly}{necessarily}
\cz{55}{response}{appeal}{outburst}{state}
\sectitle{第二节（共10小题；每小题1.5分，满分15分）}
\par 阅读下面短文，在空白处填入1个适当的单词或括号内单词的正确形式。\par
\par There's a common misconception going around that the sugar in fruit is bad for your health. So let's clear that up. The natural sugars in fruit are not the same as the sugars \blanknum{56} (add) in processed foods. In fact, fruit is incredibly \blanknum{57} (benefit) for your health. Fruits are packed with essential vitamins, minerals, fiber, water, and other nutrients \blanknum{58} are crucial for maintaining good health.\par
\par For most people, eating fruit is totally fine and encouraged \blanknum{59} part of a balanced diet. If you have diabetes or blood sugar issues, \blanknum{60} is a good idea to be mindful of portion size.\par
\par You can always try pairing your fruit with a protein source like yogurt or cheese to help \blanknum{61}(stable) blood sugar levels. The daily recommended fruit serving \blanknum{62}(vary) based on factors like age, \blanknum{63}(high), sex, weight, and physical activity. But generally, the USDA recommends that adults consume around two cups of fruit per day. What does a cup of fruit look like? 微信公众号 IAI English\par
\par Here are some \blanknum{64} (example): one small apple, one large banana, 10 dates, 22 seedless grapes, 2-3 kiwis, one large peach, three medium plums, or one slice of watermelon.\par
\par So go ahead and enjoy your favorite fruits. They're not just delicious, they're also a powerhouse of nutrition. Just remember to enjoy them in moderation and as part of a balanced diet. But please, \blanknum{65} you have chronic (慢性的) conditions that require fruit in moderation, enjoy them without concern.\par
\parttitle{第四部分 写作（共两节，满分40分）}
\sectitle{第一节（满分15分）}
\par 假设你是李华，你的学校即将举办一场国际文化节。你打算邀请你的美国朋友Jack参加。请你写一封邮件，内容包括：\par
\par (1) 活动的时间和地点；(2) 活动的目的和内容；(3) 表达你诚挚的邀请。\par
\par 注意：微信公众号 IAI English\par
\par (1) 写作词数应为 80 左右;\par
\par (2) 请按如下格式在答题纸的相应位置作答。\par
\par Dear Jack,\par
\par Yours, Li Hua\par
\sectitle{第二节（满分25分）}
\par 阅读下面材料，根据其内容和所给段落开头语续写两段，使之构成一篇完整的短文。\par
\par When I was in the fourth grade, my mom gave me a mechanical watch. It had to be rewound every night and didn't tell time as directly as an electronic watch, but it was classy and eye-catching. As one of the few people who could afford a watch, I wore it everywhere I went, subconsciously showing it off by checking the time every now and then. It left an itchy. mark stained with sweat on my wrist after being worn for a long time, but it was totally worth it when I got envious looks from peers.\par
\par In order for me to learn to manage time, my mom would casually ask me what time it was. I would roll up my sleeve, look at my watch and, after careful observation, announce the present time. Also, to make sure the time was precise, I waited patiently before the TV. When it announced the time on the hour, I meticulously adjusted the hands to ensure the watch synchronized with the TV time.\par
\par I only took off my watch before going to bed or taking a shower. It wasn't waterproof. Repairing it would cost a lot of money. My mom constantly reminded me of its value and the importance of treating it with care, saying that, once broken, it was unlikely to be replaced.\par
\par On one scorching summer day, despite parents' warnings and threats, I decided to take a dip in a lake to cool off with a bunch of other kids. Before getting in the water, I removed my watch and gently placed it near my sandals. We had a blast there, splashing water, jumping, diving, swimming and frolicking. Suddenly, we heard someone shouting, "Whose family are you kids from? How dare you? Aren't you afraid of drowning?"\par
\par 注意：微信公众号 IAI English\par
\par (1) 续写词数应为 150 左右；（2）请按如下格式在答题卡的相应位置作答。\par
\writingpara{Paragraph 1: Terrified, we scrambled up, hastily put on our clothes and fled in different directions.}
\writinglines{8}
\writingpara{Paragraph 2: I couldn't afford to let my mom know my watch was lost.}
\writinglines{8}
\end{document}
